%% kapa: exam=1 round=37 year=2026 period=2 form=A subject=accounting subject_ko=회계학 seq=37 number=77
%% source: Q-Net 감정평가사 기출문제 (cst00309) artlSeq=5258924 file=raw/1008/5258924/2026년 제37회 감정평가사 1차 2교시.pdf
%% conversion: kordoc 4.24.51 -> tools/gen1.py (auto); review=proofread-sonnet; flags=none
\begin{question}{77}
(주)감평의 20×1년도 제품 A와 B에 대한 생산ㆍ판매 자료는 다음과 같다.
\par
{\small\centering
\begin{tabular}{|l|r|r|}\hline
\multicolumn{1}{|c|}{구 분} & \multicolumn{1}{c|}{제품 A} & \multicolumn{1}{c|}{제품 B} \\\hline
예상판매수량(정규시장) & 500단위 & 300단위 \\\hline
단위당 판매가격 & ￦1,520 & ￦1,860 \\\hline
단위당 직접재료원가 & 320 & 360 \\\hline
단위당 직접노무원가 & 400 & 480 \\\hline
\makecell[l]{단위당 변동제조간접원가 및\\판매관리비} & 500 & 600 \\\hline
고정제조간접원가 및 판매관리비 & \multicolumn{2}{c|}{￦540,000} \\\hline
\end{tabular}\par}
20×1년 중 신규 고객으로부터 제품 A는 단위당 ￦1,300에 90단위, 제품 B는 단위당 ￦1,560에 110단위를 구입하겠다는 특별주문을 받았다. 연간 최대조업도는 5,000 직접노무시간이고, 직접노무시간당 임률이 ￦80일 때 특별주문을 수락할 경우 영업이익의 최대 증가액은? (단, 특별주문은 각 제품의 전부 또는 일부를 수락할 수 있으며, 기존 판매 및 원가구조에 영향을 미치지 않는다.) \par
\begin{choices}
  \item ￦7,200
  \item ￦13,200
  \item ￦13,840
  \item ￦14,200
  \item ￦14,840
\end{choices}
\end{question}
\answer{3}
